% Tema 4: conversión del PPTX original; 79 diapositivas.
% Se corrigen erratas matemáticas en signos, límites y soluciones.
\documentclass[main.tex]{subfiles}

\begin{document}
\graphicspath{{imagenes/tema4/}{../imagenes/tema4/}}


% Diapositiva original 1
\portadatema{4}{Cálculo integral. Aplicaciones (I)}

% Diapositiva original 2
\begin{frame}{Concepto de integral}
Sea $f$ continua y no negativa en $[a,b]$. Queremos calcular el área comprendida entre su gráfica y el eje $x$, desde $a$ hasta $b$.\figura{fig2.png}{.65}{.43}
\end{frame}

% Diapositiva original 3
\begin{frame}{Aproximación del área}
Podemos aproximar el área mediante polígonos. Una elección sencilla consiste en emplear rectángulos.\par\medskip Al dividir el intervalo en partes cada vez más pequeñas, obtenemos una aproximación cada vez más precisa.
\end{frame}

% Diapositiva original 4
\begin{frame}{Formalización matemática: partición}
Una partición de $[a,b]$ es un conjunto ordenado\[P=\{x_0,\ldots,x_n\},\qquad a=x_0<x_1<\cdots<x_n=b.\]Divide el intervalo en $[x_k,x_{k+1}]$, $k=0,\ldots,n-1$. Las longitudes pueden ser distintas:\[\Delta x_k=x_{k+1}-x_k.\]\figura{fig4.png}{.9}{.2}
\end{frame}

% Diapositiva original 5
\begin{frame}{Suma inferior}
En cada subintervalo tomamos\[m_k=\inf\{f(x):x\in[x_k,x_{k+1}]\}.\]Si $f$ es continua, este ínfimo es un mínimo. El área del rectángulo es $m_k(x_{k+1}-x_k)$.\figura{fig5.png}{.8}{.43}
\end{frame}

% Diapositiva original 6
\begin{frame}{Suma superior}
En cada subintervalo tomamos\[M_k=\sup\{f(x):x\in[x_k,x_{k+1}]\}.\]Si $f$ es continua, este supremo es un máximo. El área del rectángulo es $M_k(x_{k+1}-x_k)$.\figura{fig6.png}{.8}{.43}
\end{frame}

% Diapositiva original 7
\begin{frame}{Sumas inferior y superior}
\[L(f,P)=\sum_{k=0}^{n-1}m_k\Delta x_k,\qquad U(f,P)=\sum_{k=0}^{n-1}M_k\Delta x_k.\]Sumamos las áreas de todos los rectángulos.\figura{fig7.png}{.72}{.5}
\end{frame}

% Diapositiva original 8
\begin{frame}{Refinamiento de la partición}
Al refinar la partición, las sumas se aproximan al área bajo la curva. Para el límite exigimos\[\|P\|=\max_k\Delta x_k\longrightarrow0.\]\figura{fig8.png}{.72}{.5}
\end{frame}

% Diapositiva original 9
\begin{frame}{Subintervalos de igual longitud}
Si dividimos $[a,b]$ en $n$ partes iguales,\[\Delta x=\frac{b-a}{n},\qquad x_k=a+k\Delta x.\]Cuando $n\to\infty$, la anchura de los rectángulos tiende a cero.
\end{frame}

% Diapositiva original 10
\begin{frame}{Aproximación por sumas}
\[L_n=\sum_{k=0}^{n-1}m_k\Delta x,\qquad U_n=\sum_{k=0}^{n-1}M_k\Delta x.\]Para $f\ge0$ se cumple\[L_n\le\text{Área bajo la curva}\le U_n.\]Si ambos límites coinciden, proporcionan el área buscada.
\end{frame}

% Diapositiva original 11
\begin{frame}{Integral definida}
Si las sumas inferior y superior tienen el mismo límite al tender la malla a cero, la función es integrable en $[a,b]$.\par\medskip Ese límite se denomina \alert{integral definida}:\[\int_a^b f(x)\,dx.\]
\end{frame}

% Diapositiva original 12
\begin{frame}{Interpretación de la integral}
La integral representa una acumulación continua, obtenida como límite de sumas:\[\int_a^b f(x)\,dx=\lim_{\|P\|\to0}\sum_k f(\xi_k)\Delta x_k,\quad \xi_k\in[x_k,x_{k+1}].\]Para una función no negativa, es el área bajo la gráfica.
\end{frame}

% Diapositiva original 13
\begin{frame}{La integral definida es un número}
Con los límites $a$ y $b$ fijados,\[\int_a^b f(x)\,dx\]es un número. La variable $x$ es una variable de integración; puede sustituirse por otra letra.
\end{frame}

% Diapositiva original 14
\begin{frame}{Ejemplo: integral de una parábola}
Calculamos el área bajo $f(x)=x^2$ en $[0,1]$:\[\int_0^1 x^2\,dx.\]Dividimos $[0,1]$ en $n$ partes de longitud $1/n$.
\end{frame}

% Diapositiva original 15
\begin{frame}{Ejemplo: sumas inferior y superior}
\small\[L_n=\frac{1}{n^3}\sum_{k=1}^n(k-1)^2,\qquad U_n=\frac{1}{n^3}\sum_{k=1}^n k^2.\]\[\sum_{k=1}^n k^2=\frac{n(n+1)(2n+1)}6.\]\[L_n=\frac{(n-1)n(2n-1)}{6n^3},\qquad U_n=\frac{n(n+1)(2n+1)}{6n^3}.\]\figura{fig15.png}{.48}{.33}
\end{frame}

% Diapositiva original 16
\begin{frame}{Ejemplo: paso al límite}
\[\lim_{n\to\infty}L_n=\lim_{n\to\infty}U_n=\frac26=\frac13.\]Por tanto,\[\boxed{\int_0^1 x^2\,dx=\frac13.}\]\figura{fig15.png}{.5}{.35}
\end{frame}

% Diapositiva original 17
\begin{frame}{Funciones positivas y negativas}
La integral es la diferencia entre las áreas por encima y por debajo del eje $x$. El área geométrica total es $\int_a^b|f(x)|\,dx$.\[\left|\int_a^b f(x)\,dx\right|\le\int_a^b|f(x)|\,dx.\]\figura{fig17.png}{.55}{.38}
\end{frame}

% Diapositiva original 18
\begin{frame}{Propiedades de la integral definida}
\footnotesize\begin{align*}\int_a^b f&=-\int_b^a f,&\int_a^a f&=0,\\ \int_a^b kf&=k\int_a^b f,&\int_a^b(f\pm g)&=\int_a^b f\pm\int_a^b g,\\ \int_a^c f+\int_c^b f&=\int_a^b f.\end{align*}Aquí se omite $dx$ para abreviar. Para $a\le b$:\begin{itemize}\item Si $m\le f\le M$, entonces $m(b-a)\le\int_a^b f\le M(b-a)$.\item Si $f\ge g$, entonces $\int_a^b f\ge\int_a^b g$.\item Si $f\ge0$, entonces $\int_a^b f\ge0$.\end{itemize}
\end{frame}

% Diapositiva original 19
\begin{frame}{Interpretación gráfica de las propiedades}
\figura{fig19.png}{.96}{.83}
\end{frame}

% Diapositiva original 20
\begin{frame}{Teorema del valor medio integral}
Si $f$ es continua en $[a,b]$, existe $c\in[a,b]$ tal que\[\int_a^b f(x)\,dx=(b-a)f(c).\]La integral equivale al área de un rectángulo de base $b-a$ y altura igual al valor medio de la función.\figura{fig20.png}{.65}{.43}
\end{frame}

% Diapositiva original 21
\begin{frame}{Del área al cálculo de integrales}
\begin{enumerate}\item El problema del área. Concepto de integral definida.\item Teorema fundamental del cálculo integral. Regla de Barrow.\end{enumerate}
\end{frame}

% Diapositiva original 22
\begin{frame}{Primer teorema fundamental del cálculo}
Si $f$ es integrable en $[a,b]$, entonces\[F(x)=\int_a^x f(t)\,dt\]es continua en $[a,b]$.\par\medskip Si además $f$ es continua en un punto interior $c$, entonces $F$ es derivable en $c$ y\[\boxed{F\prime(c)=f(c).}\]
\end{frame}

% Diapositiva original 23
\begin{frame}{Demostración: primer teorema}
Para $f$ continua, por el teorema del valor medio,\[F(x+h)-F(x)=\int_x^{x+h}f(t)\,dt=h f(c_h),\]donde $c_h$ está entre $x$ y $x+h$. Por tanto,\[F\prime(x)=\lim_{h\to0}\frac{F(x+h)-F(x)}h=\lim_{h\to0}f(c_h)=f(x).\]\figura{fig23.png}{.9}{.25}
\end{frame}

% Diapositiva original 24
\begin{frame}{Segundo teorema fundamental del cálculo}
\small Si $f$ es continua en $[a,b]$ y $g\prime=f$, entonces\[\boxed{\int_a^b f(x)\,dx=g(b)-g(a)=[g(x)]_a^b.}\]\textbf{Demostración.} Sea $F(x)=\int_a^x f(t)\,dt$. Por el primer teorema, $F\prime=f=g\prime$, de modo que $F=g+k$.\par\medskip Como $F(a)=0$, resulta $k=-g(a)$. Así,\[F(b)=g(b)-g(a).\]
\end{frame}

% Diapositiva original 25
\begin{frame}{Primitivas y regla de Barrow}
Si $g\prime(x)=f(x)$, decimos que $g$ es una primitiva de $f$. Todas las primitivas en un intervalo son\[F(x)=g(x)+K.\]La regla de Barrow permite calcular la integral:\[\int_a^b f(x)\,dx=F(b)-F(a).\]
\end{frame}

% Diapositiva original 26
\begin{frame}{Integral indefinida}
La integral indefinida es la familia de primitivas:\[\boxed{\int f(x)\,dx=F(x)+C,\qquad F\prime=f.}\]La integración indefinida es la operación inversa de la derivación.
\end{frame}

% Diapositiva original 27
\begin{frame}{Primitivas elementales}
Para aplicar Barrow necesitamos encontrar una primitiva. Usamos tablas y reglas de integración.\begin{align*}\int0\,dx&=C,\\\int a\,dx&=ax+C\quad(a\text{ constante}),\\\int x\,dx&=\frac{x^2}{2}+C.\end{align*}
\end{frame}

% Diapositiva original 28
\begin{frame}{Potencias y logaritmos}
\begin{align*}\int x^n\,dx&=\frac{x^{n+1}}{n+1}+C,\qquad n\ne-1,\\\frac{d}{dx}\ln|x|&=\frac1x,\\\int\frac{dx}{x}&=\ln|x|+C,\\\int\frac{dx}{ax+b}&=\frac1a\ln|ax+b|+C,\qquad a\ne0.\end{align*}
\end{frame}

% Diapositiva original 29
\begin{frame}{Primitivas trigonométricas}
\[\frac{d}{dx}\sin x=\cos x,\qquad\frac{d}{dx}\cos x=-\sin x.\]\begin{align*}\int\sin x\,dx&=-\cos x+C,\\\int\cos x\,dx&=\sin x+C.\end{align*}
\end{frame}

% Diapositiva original 30
\begin{frame}{Primitiva de la exponencial}
\[\frac{d}{dx}e^x=e^x,\qquad\int e^x\,dx=e^x+C.\]\medskip Y muchas más: consultar las tablas de primitivas.
\end{frame}

% Diapositiva original 31
\begin{frame}{Integral de una función a trozos}
\small\[f(x)=\begin{cases}\cos x,&0\le x\le\pi/2,\\-1,&\pi/2<x\le\pi.\end{cases}\]\[\int_0^\pi f=\int_0^{\pi/2}\cos x\,dx+\int_{\pi/2}^{\pi}(-1)\,dx\]\[=[\sin x]_0^{\pi/2}+[-x]_{\pi/2}^{\pi}=1-\frac\pi2.\]\figura{fig31.png}{.5}{.29}
\end{frame}

% Diapositiva original 32
\begin{frame}{Integración por sustitución}
Por la regla de la cadena, si $F\prime=f$,\[\frac{d}{dx}F(g(x))=f(g(x))g\prime(x).\]Así,\[\int f(g(x))g\prime(x)\,dx=F(g(x))+C.\]Hacemos $u=g(x)$, $du=g\prime(x)\,dx$:\[\int f(u)\,du=F(u)+C.\]
\end{frame}

% Diapositiva original 33
\begin{frame}{Sustitución: ejemplo}
\[\int x\cos(x^2)\,dx.\]Tomamos $u=x^2$, $du=2x\,dx$. Entonces\[\int x\cos(x^2)\,dx=\frac12\int\cos u\,du=\frac12\sin u+C.\]\[\boxed{\int x\cos(x^2)\,dx=\frac12\sin(x^2)+C.}\]
\end{frame}

% Diapositiva original 34
\begin{frame}{Sustitución en una integral definida}
\[\int_e^5\frac{dx}{x\ln x}.\]Con $u=\ln x$, $du=dx/x$, los límites son $u(e)=1$ y $u(5)=\ln5$.\[\int_e^5\frac{dx}{x\ln x}=\int_1^{\ln5}\frac{du}{u}=[\ln u]_1^{\ln5}=\boxed{\ln(\ln5)}.\]
\end{frame}

% Diapositiva original 35
\begin{frame}{Ejercicio: sustitución}
Calcular\[\int\frac{x^2}{\sqrt{2-x}}\,dx.\]
\end{frame}

% Diapositiva original 36
\begin{frame}{Integración por partes}
\small Partimos de la derivada de un producto:\[(fg)\prime=f\prime g+fg\prime.\]Al integrar y despejar,\[\int f(x)g\prime(x)\,dx=f(x)g(x)-\int f\prime(x)g(x)\,dx.\]Escribiendo $u=f(x)$, $v=g(x)$, obtenemos\[\boxed{\int u\,dv=uv-\int v\,du.}\]
\end{frame}

% Diapositiva original 37
\begin{frame}{Cómo elegir las partes}
\[\int u\,dv=uv-\int v\,du.\]Regla mnemotécnica: «Sentado un día vi un valiente soldado vestido de uniforme».\par\medskip Para elegir $u$, suele ayudar el orden:\begin{enumerate}\item Trigonométricas inversas: $\arcsin$, $\arccos$, $\arctan$.\item Logarítmicas.\item Algebraicas, como $x^n$.\item Trigonométricas.\item Exponenciales.\end{enumerate}
\end{frame}

% Diapositiva original 38
\begin{frame}{Por partes: integral del logaritmo}
\[\int\ln x\,dx,\qquad x>0.\]Tomamos\[u=\ln x,\qquad dv=dx.\]
\end{frame}

% Diapositiva original 39
\begin{frame}{Por partes: cálculo de las partes}
\[\int\ln x\,dx.\]\[u=\ln x,\quad dv=dx,\qquad du=\frac{dx}{x},\quad v=x.\]Aplicamos $\int u\,dv=uv-\int v\,du$.
\end{frame}

% Diapositiva original 40
\begin{frame}{Por partes: resultado}
\begin{align*}\int\ln x\,dx&=x\ln x-\int x\frac{dx}{x}\\&=x\ln x-\int dx\\&=x\ln x-x+C\\&=\boxed{x(\ln x-1)+C.}\end{align*}
\end{frame}

% Diapositiva original 41
\begin{frame}{Por partes: exponencial y seno}
\[I=\int e^x\sin x\,dx.\]Tomamos $u=\sin x$, $dv=e^x dx$; entonces $du=\cos x\,dx$ y $v=e^x$.\[I=e^x\sin x-\int e^x\cos x\,dx.\]
\end{frame}

% Diapositiva original 42
\begin{frame}{Por partes: segunda aplicación}
\[I=e^x\sin x-\int e^x\cos x\,dx.\]En la nueva integral, $u=\cos x$, $dv=e^x dx$:\[\int e^x\cos x\,dx=e^x\cos x+\int e^x\sin x\,dx=e^x\cos x+I.\]Por tanto,\[I=e^x\sin x-e^x\cos x-I.\]
\end{frame}

% Diapositiva original 43
\begin{frame}{Por partes: integral cíclica}
\[I=e^x\sin x-e^x\cos x-I.\]\[2I=e^x(\sin x-\cos x).\]\[\boxed{\int e^x\sin x\,dx=\frac{e^x}{2}(\sin x-\cos x)+C.}\]
\end{frame}

% Diapositiva original 44
\begin{frame}{Aplicaciones: área entre curva y eje}
Los ceros de $f$ dividen el intervalo en regiones de signo constante. Sumamos el valor absoluto de sus integrales:\[A=\sum_i\left|\int_{c_i}^{c_{i+1}}f(x)\,dx\right|=\int_a^b|f(x)|\,dx.\]Aquí $c_0=a$ y $c_{n+1}=b$.\figura{fig44.png}{.65}{.4}
\end{frame}

% Diapositiva original 45
\begin{frame}{Área: ejemplo}
Calcular el área comprendida entre\[f(x)=x^4-5x^2+4\]y el eje $x$ en el intervalo $[-2,2]$.
\end{frame}

% Diapositiva original 46
\begin{frame}{Área: ceros y simetría}
\small Con $u=x^2$, $u^2-5u+4=0$, luego $u=4$ o $u=1$. Los ceros son $x=\pm2,\pm1$.\par\medskip La función es par, por lo que\[A=2\left(\left|\int_0^1f(x)\,dx\right|+\left|\int_1^2f(x)\,dx\right|\right).\]Una primitiva es\[F(x)=\frac{x^5}{5}-\frac{5x^3}{3}+4x.\]
\end{frame}

% Diapositiva original 47
\begin{frame}{Área: resultado}
\[F(0)=0,\qquad F(1)=\frac{38}{15},\qquad F(2)=\frac{16}{15}.\]\[A=\frac2{15}(|38|+|16-38|)=\boxed{8}.\]\figura{fig47.png}{.6}{.42}
\end{frame}

% Diapositiva original 48
\begin{frame}{Área entre una curva y ambos ejes}
Si $b$ es el primer cero pertinente de $f$ a partir de $x=0$ y no cambia de signo entre $0$ y $b$,\[A=\left|\int_0^b f(x)\,dx\right|.\]\figura{fig48.png}{.65}{.48}
\end{frame}

% Diapositiva original 49
\begin{frame}{Área con ambos ejes: ejemplo}
Calcular el área comprendida entre\[f(x)=\tan(x-1)\]y ambos ejes coordenados.
\end{frame}

% Diapositiva original 50
\begin{frame}{Área con ambos ejes: solución}
El corte con el eje $x$ es $x=1$.\begin{align*}A&=\left|\int_0^1\tan(x-1)\,dx\right|\\&=\left|[-\ln|\cos(x-1)|]_0^1\right|\\&=\boxed{-\ln(\cos1)}.\end{align*}\figura{fig48.png}{.5}{.32}
\end{frame}

% Diapositiva original 51
\begin{frame}{Ejercicio: área entre dos curvas}
Calcular el área comprendida entre\[x=y^2,\qquad y=x^2.\]\figura{fig51.png}{.7}{.55}
\end{frame}

% Diapositiva original 52
\begin{frame}{Longitud de una curva}
\small Dividimos $[a,b]$ en subintervalos. En cada uno, aproximamos la curva por un segmento:\[\Delta f_i=f(x_i)-f(x_{i-1}),\qquad \Delta x_i=x_i-x_{i-1}.\]Por Pitágoras,\[\ell_i=\sqrt{(\Delta x_i)^2+(\Delta f_i)^2}.\]\figura{fig52.png}{.6}{.38}
\end{frame}

% Diapositiva original 53
\begin{frame}{Longitud de los segmentos}
\begin{align*}\ell_i&=\sqrt{(\Delta x_i)^2+(\Delta f_i)^2}\\&=\sqrt{(\Delta x_i)^2\left[1+\left(\frac{\Delta f_i}{\Delta x_i}\right)^2\right]}\\&=\Delta x_i\sqrt{1+\left(\frac{\Delta f_i}{\Delta x_i}\right)^2}.\end{align*}
\end{frame}

% Diapositiva original 54
\begin{frame}{Aproximaciones poligonales}
Al refinar la partición, los segmentos se ajustan a la curva.\figura{fig54.png}{.98}{.58}
\end{frame}

% Diapositiva original 55
\begin{frame}{Fórmula de longitud de arco}
Para $f\in C^1([a,b])$, al pasar al límite obtenemos\begin{align*}L&=\lim_{\|P\|\to0}\sum_i\ell_i\\&=\boxed{\int_a^b\sqrt{1+[f\prime(x)]^2}\,dx.}\end{align*}
\end{frame}

% Diapositiva original 56
\begin{frame}{Longitud: ejemplo de la astroide}
Calcular la longitud de la rama superior ($y\ge0$) de\[x^{2/3}+y^{2/3}=1.\]\figura{fig56.png}{.72}{.51}
\end{frame}

% Diapositiva original 57
\begin{frame}{Astroide: derivación implícita}
\[x^{2/3}+y^{2/3}=1.\]Derivando en los puntos donde procede,\[\frac23x^{-1/3}+\frac23y^{-1/3}y\prime=0.\]\[y\prime=-\frac{y^{1/3}}{x^{1/3}}.\]\figura{fig56.png}{.5}{.29}
\end{frame}

% Diapositiva original 58
\begin{frame}{Astroide: cuadrado de la derivada}
\[(y\prime)^2=\frac{y^{2/3}}{x^{2/3}}=\frac{1-x^{2/3}}{x^{2/3}}=x^{-2/3}-1.\]Por tanto,\[\sqrt{1+(y\prime)^2}=|x|^{-1/3},\qquad x\ne0.\]\figura{fig56.png}{.55}{.35}
\end{frame}

% Diapositiva original 59
\begin{frame}{Astroide: integral de longitud}
La rama superior se extiende desde $x=-1$ hasta $x=1$.\[L=\int_{-1}^1|x|^{-1/3}\,dx=2\int_0^1x^{-1/3}\,dx.\]Es una integral impropia convergente en $x=0$.\figura{fig56.png}{.6}{.35}
\end{frame}

% Diapositiva original 60
\begin{frame}{Astroide: resultado}
\[L=2\int_0^1x^{-1/3}\,dx=2\left[\frac32x^{2/3}\right]_0^1=\boxed{3}.\]\figura{fig56.png}{.6}{.4}\footnotesize La astroide completa tiene longitud $6$.
\end{frame}

% Diapositiva original 61
\begin{frame}{Ejercicio: longitud de arco}
Calcular la longitud de la curva\[y=x^{3/2},\qquad 0\le x\le4.\]
\end{frame}

% Diapositiva original 62
\begin{frame}{Volumen por secciones transversales}
Si $A(x)$ es el área de la sección transversal de un sólido en la posición $x$, su volumen entre $a$ y $b$ es\[\boxed{V=\int_a^b A(x)\,dx.}\]\figura{fig62.png}{.65}{.45}
\end{frame}

% Diapositiva original 63
\begin{frame}{Secciones de un sólido}
\figura{fig63.png}{.85}{.75}
\end{frame}

% Diapositiva original 64
\begin{frame}{Volumen: secciones cuadradas}
Calcular el volumen de un sólido cuyas secciones transversales son cuadrados. Cada lado mide lo mismo que la cuerda vertical correspondiente de la circunferencia de radio $3$ centrada en el origen.\figura{fig64.png}{.65}{.52}
\end{frame}

% Diapositiva original 65
\begin{frame}{Volumen: solución}
\small De $x^2+y^2=9$, el lado del cuadrado es $2\sqrt{9-x^2}$.\[A(x)=4(9-x^2).\]\begin{align*}V&=2\int_0^3A(x)\,dx=8\int_0^3(9-x^2)\,dx\\&=8\left[9x-\frac{x^3}{3}\right]_0^3=\boxed{144}.\end{align*}\figura{fig64.png}{.45}{.3}
\end{frame}

% Diapositiva original 66
\begin{frame}{Ejercicio: volumen de una esfera}
Calcular, mediante secciones transversales, el volumen de una esfera de radio $r$.
\end{frame}

% Diapositiva original 67
\begin{frame}{Volumen de revolución}
Al girar alrededor del eje $x$ la región situada bajo $y=f(x)\ge0$, con $a\le x\le b$, obtenemos un sólido de revolución.\figura{fig67.png}{.9}{.55}
\end{frame}

% Diapositiva original 68
\begin{frame}{Sólidos de revolución}
La sección transversal perpendicular al eje de giro es un disco de radio $f(x)$.\figura{fig68.png}{.8}{.62}
\end{frame}

% Diapositiva original 69
\begin{frame}{Aproximación mediante discos}
\small Un disco de radio $f(\xi_i)$ y espesor $\Delta x_i$ tiene volumen\[v_i=\pi[f(\xi_i)]^2\Delta x_i.\]Sumando estos volúmenes aproximamos el volumen total.\figura{fig69.png}{.65}{.5}
\end{frame}

% Diapositiva original 70
\begin{frame}{Fórmula del volumen de revolución}
Al pasar al límite,\[V=\lim_{\|P\|\to0}\sum_i\pi[f(\xi_i)]^2\Delta x_i.\]Por tanto,\[\boxed{V=\pi\int_a^b[f(x)]^2\,dx.}\]
\end{frame}

% Diapositiva original 71
\begin{frame}{Ejemplo: volumen de un vaso}
Calcular el volumen de un vaso cuyo perfil viene dado por\[y=P(x)=\frac{2\pi}{3}x^2,\qquad 2\le y\le6.\]La altura es $y$ y el giro se realiza alrededor de ese eje; usamos la función inversa para obtener el radio.\figura{fig71.png}{.65}{.42}
\end{frame}

% Diapositiva original 72
\begin{frame}{Volumen del vaso: solución}
\small Si $t$ representa la altura, el radio es\[r(t)=\sqrt{\frac{3t}{2\pi}}.\]\[V=\pi\int_2^6r(t)^2\,dt=\int_2^6\frac{3t}{2}\,dt=\left[\frac{3t^2}{4}\right]_2^6=\boxed{24}.\]\figura{fig71.png}{.6}{.36}
\end{frame}

% Diapositiva original 73
\begin{frame}{Superficie de revolución}
La superficie lateral de un tronco de cono de radios $r$ y $R$ y generatriz $g$ es\[S=\frac{2\pi r+2\pi R}{2}\,g=\boxed{\pi(r+R)g}.\]\figura{fig73.png}{.55}{.5}
\end{frame}

% Diapositiva original 74
\begin{frame}{Troncos de cono sobre la curva}
\small En un intervalo de longitud $\Delta x$,\[r=f(x),\quad R=f(x+\Delta x),\quad\Delta f=f(x+\Delta x)-f(x).\]La generatriz es\[g=\sqrt{(\Delta x)^2+(\Delta f)^2}.\]\figura{fig74.png}{.75}{.43}
\end{frame}

% Diapositiva original 75
\begin{frame}{Suma de superficies laterales}
\small La superficie total se aproxima mediante\[S_i=\pi\bigl(f(x_{i-1})+f(x_i)\bigr)\sqrt{(\Delta x_i)^2+(\Delta f_i)^2}.\]\[S=\lim_{\|P\|\to0}\sum_i S_i.\]\figura{fig75.png}{.75}{.45}
\end{frame}

% Diapositiva original 76
\begin{frame}{Paso al límite: superficie}
\small Al tender $\Delta x$ a cero, los dos radios tienden a $f(x)$ y la generatriz aporta $\sqrt{1+[f\prime(x)]^2}\,dx$.\[S=2\pi\int_a^b f(x)\sqrt{1+[f\prime(x)]^2}\,dx.\]\figura{fig76.png}{.75}{.49}
\end{frame}

% Diapositiva original 77
\begin{frame}{Fórmula de superficie de revolución}
Para $f\ge0$, con $f\in C^1([a,b])$, la superficie generada al girar su gráfica alrededor del eje $x$ es\[\boxed{S=2\pi\int_a^b f(x)\sqrt{1+[f\prime(x)]^2}\,dx.}\]\figura{fig77.png}{.75}{.49}
\end{frame}

% Diapositiva original 78
\begin{frame}{Ejercicios}
\scriptsize\begin{enumerate}\item Calcular $\int_0^1 f(x)\,dx$, donde $0<t<1$ y\[f(x)=\begin{cases}x,&0\le x\le t,\\\displaystyle t\frac{1-x}{1-t},&t<x\le1.\end{cases}\]\item Derivar $F(x)=\displaystyle\int_{x^2}^{x^3}\frac{dt}{\sqrt{1+t^4}}$.\item Derivar $F(x)=\displaystyle\int_{\sin x}^{\cos x}\cos(\pi t^2)\,dt$.\item Calcular $\displaystyle\int_0^{2\pi}x^2\cos x\,dx$.\item Calcular $\displaystyle\int(3-x^2)^3\,dx$.\item Calcular $\displaystyle\int\frac{dx}{x+a}$, con $a$ constante.\end{enumerate}
\end{frame}

% Diapositiva original 79
\begin{frame}{Soluciones}
\footnotesize\begin{enumerate}\item $\displaystyle\int_0^1f(x)\,dx=\frac{t^2}{2}+\frac{t(1-t)}2=\boxed{\frac t2}$.\item $\displaystyle F\prime(x)=\frac{3x^2}{\sqrt{1+x^{12}}}-\frac{2x}{\sqrt{1+x^8}}$.\item $\displaystyle F\prime(x)=-\sin x\cos(\pi\cos^2x)-\cos x\cos(\pi\sin^2x)$\\$\displaystyle\phantom{F\prime(x)}=(\sin x-\cos x)\cos(\pi\sin^2x)$.\item $\displaystyle\int_0^{2\pi}x^2\cos x\,dx=4\pi$.\item $\displaystyle\int(3-x^2)^3\,dx=27x-9x^3+\frac95x^5-\frac17x^7+C$.\item $\displaystyle\int\frac{dx}{x+a}=\ln|x+a|+C$.\end{enumerate}
\end{frame}
\end{document}
